\documentclass[12pt]{article}

\usepackage{amsmath, amssymb, amsthm}
\usepackage{geometry}
\usepackage{hyperref}
\usepackage{microtype}
\usepackage{parskip}
\usepackage{titlesec}
\usepackage{abstract}
\usepackage{booktabs}
\usepackage{mathtools}

\geometry{
  paper=letterpaper,
  top=1.25in,
  bottom=1.25in,
  left=1.25in,
  right=1.25in
}

\hypersetup{
  colorlinks=true,
  linkcolor=black,
  citecolor=black,
  urlcolor=black
}

\titleformat{\section}{\large\bfseries}{\thesection.}{0.75em}{}
\titleformat{\subsection}{\normalsize\bfseries}{\thesubsection.}{0.75em}{}

\newtheorem{definition}{Definition}
\newtheorem{theorem}{Theorem}
\newtheorem{proposition}{Proposition}

\title{\textbf{Trajectory and Residue}\\[0.5em]
\large How Work Becomes What Continues to Hold}
\author{Flyxion}
\date{}

\begin{document}

\maketitle

\begin{abstract}
We formalize a selection-theoretic account of why institutions
systematically reward signal management over real work, and propose an
alternative infrastructure in which work is evaluated and compensated
according to what persists under pressure. Modern institutions optimize
for the management of signals rather than the closure of constraints.
This produces a stable class of
roles---identified by David Graeber as ``bullshit jobs'' \cite{graeber2018}---whose
function is not to resolve real-world pressures but to mediate,
translate, and amplify representations of work. This document
formalizes that diagnosis as a selection problem and proposes an
alternative infrastructure in which work is evaluated, routed, and
compensated according to persistent constraint closure. The core shift
replaces activity-based metrics with trajectory-based residue: evidence
that a system resisted, was repaired, and remained stable over time.
Three adversarial refinements---causal dependency, distributed
attribution, and stress-weighted persistence---are incorporated to
ensure the residue function cannot be gamed without genuine interaction
with real constraints. The result is a framework for reorganizing labor
markets around verifiable interaction with reality rather than
credentialed coordination.
\end{abstract}

\section{The Existing Selection Function}

Let a system of work allocation be defined by a selection function
$\mathcal{S}$ that maps observed signals $\sigma$ to resource allocation
decisions $R$:
\[
  \mathcal{S} : \sigma \rightarrow R
\]
In current institutional environments, $\sigma$ consists primarily of:
\begin{itemize}
  \item activity counts (tickets, commits, meetings);
  \item credentials (degrees, titles, affiliations);
  \item narrative coherence (reports, presentations, alignment language).
\end{itemize}

These signals are not required to correspond to constraint closure.
Instead, they are optimized for legibility to management structures.
The resulting equilibrium is characterized by high signal production,
low constraint interaction, and stable mediation layers.

This is the formal condition for the persistence of bullshit jobs \cite{graeber2018}: roles
that maximize $\sigma$ without engaging the underlying system
constraints. When a measure becomes a target, it ceases to be a good
measure \cite{goodhart1975,strathern1997}.

\begin{proposition}
Under $\mathcal{S}$, the dominant strategy for an agent seeking to
maximize $R$ is to produce high-volume, low-cost signals rather than to
close constraints. Since signals are cheaper to produce and more directly
legible to $\mathcal{S}$ than constraint closure, agents rationally
allocate effort toward signal production \cite{simon1955,arrow1970}.
Signal production and constraint closure are therefore decoupled at
equilibrium.
\end{proposition}

The managerial, credentialing, and recruiter layers that characterize
contemporary labor markets are not inefficiencies---they are stable
equilibria of $\mathcal{S}$ \cite{lucas1976}.

\section{Constraint Closure as the Primitive}

\begin{definition}[Constraint]
A \emph{constraint} is a condition imposed by reality on a system's
behavior. A constraint $c \in C$ is \emph{closed} when the system
reaches a state in which $c$ is satisfied and remains satisfied under
perturbation.
\end{definition}

Let $C$ be a set of constraints over a system $X$. A closure event is a
transition:
\[
  X_t \rightarrow X_{t+1} \quad \text{such that} \quad C(X_{t+1}) = 0
\]
where $C(\cdot) = 0$ indicates satisfaction.

Instantaneous satisfaction is insufficient. A constraint closed at
$t+1$ may reopen under load, under compositional stress, or as
downstream systems evolve. Closure is therefore not a point event but a
trajectory-dependent quantity, whose value is determined by what
survives subsequent interaction. The operative notion is therefore not
satisfaction but \emph{residue}: the durable trace left by a closure
event in a system subject to ongoing interaction.

\section{Residue: The Basic Definition}

\begin{definition}[Closure Residue]
Let $r$ be a repair event that closes a constraint $c$ at time $t$.
The \emph{closure residue} of $r$ over window $[t, t+\Delta]$ is:
\[
  \rho(r) = \int_{t}^{t+\Delta}
    \mathbb{I}\bigl[C(X_\tau) = 0\bigr]
    \cdot D(\tau)
    \cdot S(\tau)
  \, d\tau
\]
where $D(\tau)$ is the causal dependency weight at time $\tau$ and
$S(\tau)$ is the stress intensity at time $\tau$.
\end{definition}

Any definition of residue that can be satisfied without interacting with
real constraints will be optimized for without producing value. The
three terms above are therefore not design choices but adversarial
requirements: each must be defined so that gaming it is equivalent to
doing the work. The following three sections address them in turn.

\section{Causal Dependency}

The naive reading of $D(\tau)$ is structural: system $B$ \emph{declares}
that it depends on system $A$. This is exploitable. Agents can inflate
dependency graphs to amplify residue scores without producing causal
reliance.

The correct formulation treats dependency as \emph{sensitivity under
perturbation}. System $B$ causally depends on closure $r$ if and only
if the removal of $r$ increases the probability of failure in $B$.

\begin{definition}[Causal Dependency Weight]
\[
  D(\tau) = \sum_{j \in \mathrm{downstream}}
    w_j \cdot \Delta P\!\left(\mathrm{failure}_j \mid \neg r,\, \tau\right)
\]
where $\Delta P(\mathrm{failure}_j \mid \neg r, \tau)$ is the increase
in failure probability for downstream system $j$ when the repair $r$ is
counterfactually removed at time $\tau$, and $w_j$ is a weight
proportional to the criticality of system $j$.
\end{definition}

This formulation aligns the framework with interventionist causality
\cite{pearl2009,spirtes2000}.
Dependency is not declared; it is measured by failure propagation. This
differs from structural dependency graphs, which encode declared
relationships rather than causal sensitivity under intervention. A
repair that no downstream system relies on---regardless of how many
systems annotate it as a dependency---contributes zero to $D(\tau)$.

\begin{definition}[Metric vs.\ Constraint]
A \emph{metric} $m : X \to \mathbb{R}$ is an observation channel. A
\emph{constraint} $c(X, E) = 0$ is a resistance relation whose
satisfaction depends on interaction with an environment $E$ that the
evaluated agent does not fully control. Metrics can be fabricated by
manipulating what the system records. Constraints cannot be fabricated
without changing what the system must survive.
\end{definition}

This distinction is load-bearing. Any residue term reducible to a metric
controlled by the evaluated agent collapses back into $\mathcal{S}$.

\section{Distributed Attribution}

The current draft assumes that a repair $r$ can be associated cleanly
with a single agent. In practice, most nontrivial closures are
composite: partial fixes, prior groundwork, review interventions, and
subsequent stabilizations each contribute to the final residue.

Without explicit handling, the system drifts toward last-touch
attribution: the agent who closes the ticket receives full credit,
regardless of who performed the diagnostic work, who identified the root
cause, or who stabilized the environment that made repair possible. This
is easy to exploit.

The correct approach treats each closure event as a directed acyclic
graph (DAG) of contributing actions, and assigns credit via a
marginal-contribution decomposition analogous to cooperative game
attribution \cite{shapley1953,roughgarden2016}.

\begin{definition}[Marginal Contribution]
Let $\mathcal{A} = \{a_1, \ldots, a_n\}$ be the set of actions
contributing to closure residue $\rho$. The marginal contribution of
action $a_i$ is:
\[
  \rho_i =
    \mathbb{E}\bigl[\rho \mid \text{with } a_i\bigr]
    - \mathbb{E}\bigl[\rho \mid \text{without } a_i\bigr]
\]
\end{definition}

This is a Shapley-style decomposition \cite{shapley1953} over the repair graph. It is
computationally expensive in the general case, but approximate versions
operating over observable repair DAGs would substantially improve
resistance to gaming. The key property is that credit is \emph{marginal
contribution to closure}, not proximity to the final commit or the final
ticket resolution.

Agents who perform diagnostic groundwork, stabilize test environments,
or review repairs that would otherwise have introduced regressions
accumulate residue proportional to their counterfactual impact.
Agents who append low-value actions to high-residue repair chains do
not.

\section{Stress-Weighted Persistence}

Duration alone is insufficient as a persistence measure. A trivial fix
that is never exercised again will appear persistent simply because
nothing has reopened the constraint. What matters is not whether a
closure holds over time, but whether it holds \emph{under conditions
that would reveal failure}.

\begin{definition}[Stress Term]
$S(\tau)$ measures the intensity of system interaction, load, or
adversarial conditions at time $\tau$. Operationally, it can be
estimated from:
\begin{itemize}
  \item throughput and request volume;
  \item the frequency of adjacent constraint activations;
  \item the presence of concurrent incidents in dependent systems.
\end{itemize}
\end{definition}

With the stress term incorporated, the residue integral rewards closures
that survive load. A repair holding under high $S(\tau)$ accumulates
residue rapidly. A repair in a system that has gone dormant accumulates
little, regardless of how long it nominally persists. In software
systems, this corresponds to fixes that hold under production traffic
rather than only under test conditions.

This prevents dead zones---abandoned subsystems, deprecated
codepaths---from generating artificial residue through inactivity.

\section{The Refined Residue Function}

Assembling the three refinements:

\begin{definition}[Refined Closure Residue]
\[
  \rho(r) = \int_{t}^{t+\Delta}
    \mathbb{I}\bigl[C(X_\tau) = 0\bigr]
    \cdot
    \left(
      \sum_{j \in \mathrm{downstream}}
        w_j \cdot \Delta P\!\left(\mathrm{failure}_j \mid \neg r,\, \tau\right)
    \right)
    \cdot S(\tau)
  \, d\tau
\]
with marginal attribution $\rho_i =
\mathbb{E}[\rho \mid \text{with } a_i] -
\mathbb{E}[\rho \mid \text{without } a_i]$
distributed across the contributing action DAG.
\end{definition}

This function has the following adversarial properties:
\begin{itemize}
  \item \textbf{Dependency must be causal.} Annotating connections that
    do not correspond to actual failure propagation does not increase
    $D(\tau)$.
  \item \textbf{Attribution is distributed.} Last-touch exploitation is
    suppressed; marginal contribution over the full repair DAG
    determines credit allocation.
  \item \textbf{Persistence must be exercised.} Closures accumulate
    residue only under conditions of real system stress. Dormant fixes
    do not.
\end{itemize}

Faking high residue under this function requires actually interacting
with real constraints under real load. That is not gaming the
system---it is doing the work.

\subsection{Exogeneity Condition}

A further adversarial pressure is synthetic stress: an agent controls
both the intervention and the environment that validates it, generating
artificial load to inflate $S(\tau)$. Bot traffic, fake dependencies,
synthetic incidents, and self-generated load must therefore be
discounted unless they correspond to externally consequential use.

The requirement is that valid stress be causally tied to real downstream
reliance, not merely to volume. Define:

\begin{definition}[Valid Stress]
\[
  S_{\mathrm{valid}}(\tau) = S(\tau) \cdot E_{\mathrm{exo}}(\tau) \cdot D(\tau)
\]
where $E_{\mathrm{exo}}(\tau) \in [0,1]$ measures the independence of
the stress source from the agent being evaluated. Stress generated or
controlled by the evaluated agent is discounted toward zero.
\end{definition}

Stress counts only when paired with causal dependency and external
consequence. The exogeneity condition prevents an agent from being both
the source and the beneficiary of the load that validates their work.

\subsection{Maintenance Residue}

The repair-event formulation risks a perverse incentive: agents may
allow constraints to fail in order to generate repair events and
accumulate residue. The fix is to distinguish repair residue from
maintenance residue.

Repair residue $\rho_{\mathrm{repair}}(r)$ arises from closing a
violated constraint. Maintenance residue arises from keeping a
constraint closed under valid stress in the absence of a discrete
failure event:

\begin{definition}[Maintenance Residue]
\[
  \rho_{\mathrm{maint}}(a) =
    \int_{t}^{t+\Delta}
      \mathbb{I}\bigl[C(X_\tau) = 0\bigr]
      \cdot D(\tau)
      \cdot S_{\mathrm{valid}}(\tau)
      \cdot M_a(\tau)
    \, d\tau
\]
where $M_a(\tau)$ estimates the marginal contribution of agent $a$'s
preventative actions to continued closure at time $\tau$.
\end{definition}

Total residue is then:
\[
  \rho_{\mathrm{total}} = \rho_{\mathrm{repair}} + \rho_{\mathrm{maint}}
\]

This rewards the agent whose work makes failure less likely in the first
place, not only the agent who resolves failures after they occur.

\section{Trajectories Rather Than Snapshots}

Let a worker's contribution be modeled as a trajectory through
constraint space:
\[
  \gamma : t \mapsto X_t
\]
The value of the trajectory is not the number of interventions but the
cumulative residue of closures along it, weighted toward recent
persistence under stress:
\[
  V(\gamma) = \sum_i \rho_i \cdot e^{-\lambda(t - t_i)}
\]
where $\lambda > 0$ is a decay parameter. This prevents static
reputations from dominating: residue from closures that are no longer
stress-tested decays, while active maintenance under load continues to
accumulate. This replaces pointwise metrics with path-dependent
evaluation.

Two trajectories with equal activity counts can differ dramatically:

\medskip
\begin{center}
\begin{tabular}{lcc}
\toprule
Trajectory type & Activity & $V(\gamma)$ \\
\midrule
High activity, low persistence & High & Low \\
Low activity, high persistence under stress & Low & High \\
\bottomrule
\end{tabular}
\end{center}
\medskip

The current selection function $\mathcal{S}$ cannot distinguish these
trajectories. The replacement function $\mathcal{S}'$ can.

\section{The Replacement Selection Function}

\begin{definition}[Residue-Based Selection Function]
\[
  \mathcal{S}' : \rho(\gamma) \rightarrow R
\]
Resources---tasks, compensation, authority---are allocated according to
accumulated and recent closure residue, distributed via marginal
attribution over contributing action DAGs.
\end{definition}

The equilibrium under $\mathcal{S}'$ differs structurally from the
equilibrium under $\mathcal{S}$: only actions that produce persistent,
stress-tested, causally integrated closure accumulate value; all other
signals are neutral or decaying.

Bullshit jobs---roles that maximize $\sigma$ under $\mathcal{S}$---do
not satisfy $\mathcal{S}'$. They are not filtered by gatekeeping; they
simply fail to accumulate residue.

The distinction between $\mathcal{S}$ and $\mathcal{S}'$ is not
normative but structural: one operates on signals independent of
constraint interaction; the other operates on residue produced by it.

\section{Economic Routing}

Measurement alone is insufficient unless it controls allocation.
$\mathcal{S}'$ must route not just evaluation but resources: work
opportunities, payment, and decision authority.

\medskip
\noindent\textit{Without economic routing, the residue framework becomes
a descriptive overlay---a better dashboard for the same managerial layer
it is intended to replace.}
\medskip

The mechanism must place resource allocation decisions downstream of
$\rho(\gamma)$, bypassing institutional hierarchy.

The near-term viable domain for this is:
\begin{itemize}
  \item contract and freelance markets where buyers can observe
    constraint surfaces directly;
  \item open systems where incident and dependency graphs are partially
    public;
  \item environments where the cost of failure is immediate and
    attributable.
\end{itemize}

In these domains, buyers can select contributors based on demonstrated
closure trajectories without requiring credential intermediaries or HR
gatekeeping \cite{coase1937,williamson1985}. The recruiter and
credentialing layers do not get reformed; they get bypassed.

\section{Structural Resistance}

Bullshit jobs persist because they serve real interests. They provide
control surfaces for institutional hierarchy. They stabilize internal
power structures by converting uncertain, judgment-dependent work into
predictable, legible signals. They give institutions a way to manage
headcount without depending on autonomous technical judgment.

A residue-based system therefore does not merely compete technically
with existing infrastructure. It competes with existing interests.

Adoption will occur where:
\begin{itemize}
  \item constraint visibility is high and incident surfaces are
    observable;
  \item the cost of failure is immediate and cannot be absorbed by
    mediation;
  \item mediation layers are already thin.
\end{itemize}

Adoption will be resisted where signaling is more valuable than
outcomes---which describes most large institutional environments.

This is not an argument against building the system. It is a prediction
about where it will first take hold and where it will face the most
friction.

\section{Implementation Sketch}

A minimal viable implementation does not attempt full generality. It
chooses a domain where the incident surface is observable, the
dependency structure is partially inferable, and stress signals exist.

The system requires:
\begin{enumerate}
  \item Data ingestion from observable systems (repositories, logs,
    incident trackers).
  \item Construction of incident and dependency graphs.
  \item Attribution of repairs to contributing action DAGs.
  \item Temporal tracking of persistence under stress.
  \item Computation of residue scores via marginal attribution.
  \item Interface for querying contributor trajectories.
  \item Mechanism for routing tasks and payment based on accumulated
    residue.
\end{enumerate}

The validation criterion is discriminative power under equal activity:
given two contributors with similar activity levels, does the
residue-based metric separate them in a way that aligns with expert
judgment? If that holds even weakly, the architecture becomes
economically interesting and the selection-function argument moves from
structural to demonstrable.

\section{Evaluation Without a Constraint Surface}

The framework developed above applies directly to labor markets, where
constraint surfaces are partially observable and residue can be
approximated from incident and dependency data. But the same
selection-theoretic logic applies wherever evaluation has become
decoupled from the underlying structure it once tracked.

Consider a system in which agents historically accumulated residue
through repeated interaction with stable constraint sets embedded in
social roles. The role provided the trajectory; the community provided
the measurement; observable success and failure under stress provided
the residue. Evaluation vocabulary---the words a system uses to
distinguish trajectories---was compressed residue: lossy encodings of
long-run closure patterns that allowed agents to communicate about
trajectory differences without enumerating them.

When the institutional structures that exposed those constraint sets
dissolve, the collapse is structural and forced. If trajectories cannot
be evaluated against a shared constraint surface, then residue cannot be
accumulated or compared. If residue cannot be accumulated or compared,
evaluation vocabulary ceases to map to observable distinctions. When
vocabulary ceases to map to observable distinctions, it decays.
Evaluation then defaults to signals that remain legible under
$\mathcal{S}$: appearance, affiliation, stated identity.

\medskip
\noindent\textit{This is not primarily a cultural or emotional
phenomenon. It is a measurement collapse: the system can no longer
distinguish trajectories that differ in constraint closure, because the
constraint surface is no longer exposed.}
\medskip

\section{Resentment as Equilibrium Behavior}

Under $\mathcal{S}$, agents who cannot produce high residue but are
still ranked by visible signals face a structural problem with no
internal solution. Under $\mathcal{S}$ without access to observable
$\rho$, no strategy based on deeper constraint interaction yields higher
expected $R$. The only remaining degrees of freedom are in the
interpretation and production of $\sigma$.

The dominant response is therefore not to interact more deeply with
constraints---that path is closed---but to revalue the ranking system
itself:
\[
  \text{inaccessible } \rho
  \;\Rightarrow\;
  \text{evaluation collapse}
  \;\Rightarrow\;
  \text{signal inversion}
\]

This is the structural mechanism behind what is sometimes described as
resentment or transvaluation \cite{kahneman1979}: not primarily an
emotional state, but the optimal adaptation to operating under
$\mathcal{S}$ without access to $\mathcal{S}'$.

The behavior is stable as long as the measurement layer remains
collapsed. Reintroducing observable constraint surfaces---roles,
systems, environments where closure is visible and stress-testable---
changes the equilibrium without requiring any change in the agents
themselves. The problem is not the agents. The problem is the missing
measurement layer.

\section{Role Dissolution as Trajectory Collapse}

Historically, stable social roles functioned as constrained trajectory
spaces. A role instantiated a persistent mapping $\gamma$ under a fixed
constraint set $C$, with publicly observable closure and failure events.
The role was not merely a job description---it was a scaffold for
accumulating residue in a shared measurement system. Evaluation
vocabulary was calibrated against that accumulation.

When roles dissolve, the shared constraint set $C$ and the
observability of closure events are removed simultaneously, eliminating
the possibility of shared residue accumulation. Agents are left in
unconstrained space with no stable external basis for evaluation. The
vocabulary that once described trajectory differences becomes decorative
rather than functional, because the underlying structures it referred to
are no longer exposed.

The repair is not to restore the old roles. It is to restore exposure
to constraint surfaces: environments where systems resist, repairs are
visible, and closures persist under stress. The vocabulary re-emerges
after the measurement system is restored, not before.

This is a direct extension of $\mathcal{S}'$ beyond labor markets. The
replacement selection function does not require shared tradition,
inherited language, or institutional continuity. It requires only that
constraint surfaces be observable, that closures be attributable, and
that residue accumulate in proportion to genuine interaction with
reality.

\section{Dynamic Constraint Fields}

The preceding formulation treats constraints as discrete conditions
evaluated over system states. In practice, constraints evolve as
functions of both system configuration and environmental interaction.
We therefore refine the model by treating constraints as a field over
state space.

Let the system be defined over a domain $\Omega$ with state
$X : \Omega \times \mathbb{R} \rightarrow \mathcal{X}$. A constraint
field is a mapping
\[
  \mathcal{C} : \mathcal{X} \times \mathcal{E} \rightarrow \mathbb{R}^k
\]
such that $\mathcal{C}(X,E)=0$ defines the manifold of satisfied
constraints.

Closure is no longer a pointwise event but a trajectory that remains
within the zero-set of $\mathcal{C}$ under system evolution:
\[
  \forall \tau \in [t,t+\Delta], \quad
  \mathcal{C}(X_\tau, E_\tau) = 0.
\]

This formulation allows constraints to shift, tighten, or couple over
time \cite{kauffman1993}. Residue therefore measures not only persistence
but stability under a moving constraint surface. The residue integral remains
well-defined under this generalization, with the indicator function
replaced by membership in the constraint manifold.

\section{Adversarial Behavior and Gaming Resistance}

Any selection function that allocates resources induces strategic
behavior. A valid framework must therefore be evaluated not only under
honest participation but under adversarial optimization.

Let an agent choose actions $a \in \mathcal{A}$ to maximize expected
reward under $\mathcal{S}'$. A gaming strategy is any policy $\pi$
that increases $\mathbb{E}[V(\gamma)]$ without increasing true
constraint closure.

The primary attack vectors are as follows.

\medskip
\noindent\textbf{Synthetic Dependency Inflation.}
Agents construct artificial dependency graphs to increase $D(\tau)$.
This is neutralized by counterfactual dependency measurement:
non-causal edges contribute zero.

\medskip
\noindent\textbf{Stress Fabrication.}
Agents generate load to increase $S(\tau)$.
This is neutralized by the exogeneity filter $E_{\mathrm{exo}}(\tau)$,
which discounts self-generated stress.

\medskip
\noindent\textbf{Attribution Capture.}
Agents position themselves at terminal points of repair chains to
capture credit. This is neutralized by marginal contribution over the
repair DAG.

\medskip

\begin{proposition}
Under $\mathcal{S}'$, any strategy that increases expected reward must
increase expected interaction with exogenous constraint surfaces under
stress.
\end{proposition}

\begin{proof}[Sketch]
Each multiplicative term in $\rho$ is independently grounded in
counterfactual or exogenous measurement. Artificial inflation of any
single term is nullified unless it produces corresponding changes in
failure propagation or externally induced stress. Therefore, increasing
$\rho$ requires genuine causal contribution to constraint closure.
\end{proof}

\section{Equilibrium Analysis}

We compare the equilibria induced by $\mathcal{S}$ and $\mathcal{S}'$
in a simplified agent model \cite{osborne1994,aumann1974}.

Let agents choose between two action classes: signal production
$a_\sigma$ and constraint interaction $a_c$. Let costs satisfy
$C_\sigma \ll C_c$, and let rewards depend on the selection function.

Under $\mathcal{S}$:
\[
  R_\sigma \gg R_c,
\]
since signals are directly observed and rewarded.

Under $\mathcal{S}'$:
\[
  R_\sigma \approx 0, \qquad R_c \propto \rho.
\]

\begin{proposition}
The Nash equilibrium under $\mathcal{S}$ is dominated by signal
production. The Nash equilibrium under $\mathcal{S}'$ is dominated by
constraint interaction.
\end{proposition}

\begin{proof}[Sketch]
Under $\mathcal{S}$, the payoff-to-cost ratio of $a_\sigma$ exceeds that
of $a_c$, so all agents shift toward signal production. Under
$\mathcal{S}'$, $a_\sigma$ yields negligible reward, while $a_c$
produces residue proportional to constraint closure. The payoff
dominance reverses, and agents converge to constraint interaction.
\end{proof}

The transition from $\mathcal{S}$ to $\mathcal{S}'$ therefore induces a
phase shift in agent behavior, not a marginal adjustment.

\section{Information-Theoretic Interpretation}

Residue can be interpreted as a measure of irreversible information
imprinted into a system through constraint interaction
\cite{landauer1961,bennett1982,lloyd2000}.

Let $H(X)$ denote the entropy of system state $X$. A closure event that
stabilizes a constraint reduces the accessible state space of the
system. Under repeated stress, this reduction becomes persistent.

\begin{definition}[Informational Residue]
\[
  \rho_{\mathrm{info}}(r)
  =
  \int_{t}^{t+\Delta}
    \Delta H_{\mathrm{irreversible}}(\tau)
    \cdot D(\tau)
    \cdot S_{\mathrm{valid}}(\tau)
  \, d\tau
\]
where $\Delta H_{\mathrm{irreversible}}(\tau)$ is the irreversible
reduction in system entropy induced by the closure at time $\tau$.
\end{definition}

This formulation connects residue to thermodynamic work
\cite{georgescu1971,bennett1982}. A closure that
does not persist corresponds to reversible fluctuation and leaves no
lasting informational trace. Only closures that reduce entropy under
stress contribute to residue.

This interpretation clarifies why signal production fails to substitute
for constraint interaction: signals can be generated without altering
the entropy structure of the underlying system.

\section{Interface and Observability}

The effectiveness of $\mathcal{S}'$ depends on the observability of
constraint interaction. Let $\mathcal{O}$ be the observation operator
mapping system dynamics to measurable data:
\[
  \mathcal{O} : (X_t, E_t) \rightarrow \text{logs, traces, events.}
\]

A necessary condition for implementing $\mathcal{S}'$ is that
$\mathcal{O}$ captures constraint activation and resolution events,
failure propagation across system components \cite{barabasi1999,newman2010},
stress signals independent of agent control, and temporal persistence
under load.

Incomplete observability introduces bias into $\rho$ but does not
revert the system to signal-based selection. Instead, it degrades
measurement resolution. The design problem is therefore not to eliminate
measurement error but to ensure that errors are orthogonal to strategic
manipulation. As long as measurement noise is not controllable by
agents, the adversarial properties of $\mathcal{S}'$ are preserved.

\section{Conclusion}

Graeber \cite{graeber2018} identified a class of roles that optimize for
signal management rather than real work. Their persistence is not accidental. It is the
stable equilibrium of a selection function $\mathcal{S}$ that rewards
legibility over constraint interaction.

Replacing that function requires three moves: redefining work in terms
of trajectory and residue; constructing measurement systems that capture
causal dependency, distributed attribution, and stress-weighted
persistence; and routing economic decisions through those measurements
rather than through credential and title hierarchy.

The goal is not to eliminate coordination. Coordination that closes
constraints accumulates residue. The goal is to subordinate coordination
to constraint closure, so that roles which exist only to manage
representations of work can no longer satisfy the selection function
that allocates resources.

\medskip
\noindent The next step is a toy model: agents choosing between signal
production and constraint interaction under $\mathcal{S}$ and
$\mathcal{S}'$, with equilibria computed explicitly for each. That model
would make the argument not just structural but demonstrable.

\medskip
\begin{center}
\textit{Work becomes not what is reported, but what continues to hold.}
\end{center}

\begin{thebibliography}{99}

\bibitem{graeber2018}
D.~Graeber,
\textit{Bullshit Jobs: A Theory}.
Simon \& Schuster, 2018.

\bibitem{goodhart1975}
C.~A.~E.~Goodhart,
``Problems of Monetary Management: The U.K. Experience,''
in \textit{Papers in Monetary Economics}, Reserve Bank of Australia, 1975.

\bibitem{strathern1997}
M.~Strathern,
``Improving Ratings: Audit in the British University System,''
\textit{European Review}, vol.~5, no.~3, pp.~305--321, 1997.

\bibitem{lucas1976}
R.~E.~Lucas,
``Econometric Policy Evaluation: A Critique,''
\textit{Carnegie-Rochester Conference Series on Public Policy}, 1976.

\bibitem{pearl2009}
J.~Pearl,
\textit{Causality: Models, Reasoning, and Inference}, 2nd~ed.
Cambridge University Press, 2009.

\bibitem{spirtes2000}
P.~Spirtes, C.~Glymour, and R.~Scheines,
\textit{Causation, Prediction, and Search}.
MIT Press, 2000.

\bibitem{shapley1953}
L.~S.~Shapley,
``A Value for $n$-Person Games,''
in \textit{Contributions to the Theory of Games II},
Princeton University Press, 1953.

\bibitem{roughgarden2016}
T.~Roughgarden,
\textit{Twenty Lectures on Algorithmic Game Theory}.
Cambridge University Press, 2016.

\bibitem{osborne1994}
M.~J.~Osborne and A.~Rubinstein,
\textit{A Course in Game Theory}.
MIT Press, 1994.

\bibitem{arrow1970}
K.~J.~Arrow,
\textit{Essays in the Theory of Risk-Bearing}.
North-Holland, 1970.

\bibitem{coase1937}
R.~H.~Coase,
``The Nature of the Firm,''
\textit{Economica}, vol.~4, no.~16, pp.~386--405, 1937.

\bibitem{williamson1985}
O.~E.~Williamson,
\textit{The Economic Institutions of Capitalism}.
Free Press, 1985.

\bibitem{landauer1961}
R.~Landauer,
``Irreversibility and Heat Generation in the Computing Process,''
\textit{IBM Journal of Research and Development}, vol.~5, no.~3, 1961.

\bibitem{bennett1982}
C.~H.~Bennett,
``The Thermodynamics of Computation---A Review,''
\textit{International Journal of Theoretical Physics}, 1982.

\bibitem{lloyd2000}
S.~Lloyd,
``Ultimate Physical Limits to Computation,''
\textit{Nature}, vol.~406, pp.~1047--1054, 2000.

\bibitem{kahneman1979}
D.~Kahneman and A.~Tversky,
``Prospect Theory: An Analysis of Decision under Risk,''
\textit{Econometrica}, vol.~47, no.~2, 1979.

\bibitem{simon1955}
H.~A.~Simon,
``A Behavioral Model of Rational Choice,''
\textit{Quarterly Journal of Economics}, 1955.

\bibitem{aumann1974}
R.~J.~Aumann,
``Subjectivity and Correlation in Randomized Strategies,''
\textit{Journal of Mathematical Economics}, 1974.

\bibitem{barabasi1999}
A.-L.~Barab\'{a}si and R.~Albert,
``Emergence of Scaling in Random Networks,''
\textit{Science}, 1999.

\bibitem{newman2010}
M.~E.~J.~Newman,
\textit{Networks: An Introduction}.
Oxford University Press, 2010.

\bibitem{kauffman1993}
S.~A.~Kauffman,
\textit{The Origins of Order: Self-Organization and Selection in Evolution}.
Oxford University Press, 1993.

\bibitem{georgescu1971}
N.~Georgescu-Roegen,
\textit{The Entropy Law and the Economic Process}.
Harvard University Press, 1971.

\end{thebibliography}

\end{document}
